\subsection{Prusti}

Prusti is a verification framework whose development has started in 2017 at the ETH Zurich \cite{astrauskas_leveraging_2019,deshmukh_prusti_2022}. It leverages Rust's type system while providing a specification language, which can be used to describe certain properties of the program for verifying \cite{astrauskas_leveraging_2019}. This verifying process is automated by translating the Rust program and specifications into the Viper \cite{jobstmann_viper_2016} intermediate language, which is used for reasoning \cite{astrauskas_leveraging_2019}. There were three main goals while designing Prusti:

\begin{itemize}
    \item Provide functional correctness properties for Rust programs, without depending on Rust's safety guarantees \cite{deshmukh_prusti_2022}.
    \item Make verification accessible to programmers without requiring the domain knowledge of verification \cite{deshmukh_prusti_2022}.
    \item Make the usage of Prusti easy for Rust programmers \cite{deshmukh_prusti_2022}.
\end{itemize}

\subsubsection{Usability}
\label{cap:prusti_usability}

Usability is prioritized goal for Prusti.
Prusti integrates itself into the Rust compiler and uses the generated \acs{mir} for further interaction \cite{deshmukh_prusti_2022}. This has the advantage that the verifier itself always stays on pair with the progress of the Rust language, without relying on adaptions of the tool \cite{deshmukh_prusti_2022}. Another advantage is the resulting ease of use with \textit{cargo}, which allows Prusti a convenient integration in an already existing toolchain for Rust development. Furthermore, there exists an extension called Prusti Assistant \cite{chair_of_programming_methodology_-_eth_zurich_prusti_nodate} for the popular code editor \textit{Visual Studio Code}. It takes care of the local setup process of Prusti and integrates itself into the editor for inline error highlighting of error messages produced by Prusti.

Prusti is fully autonomous. The MIR, generated by the compiler, gets translated \cite{astrauskas_leveraging_2019} into the Viper intermediate language, which is a part of the Viper infrastructure \cite{jobstmann_viper_2016}. The Viper infrastructure then produces the verification results and reports them back to Prusti, which then consequently reports them back to the user \cite{astrauskas_leveraging_2019}.

Prusti proposes incremental verification. The tool makes it easy to retrofit an existing codebase with Prusti annotations. Prusti proves by default that a program will not terminate after reaching an unrecoverable state, e.g. a panic!() or unreachable!() macro call, without the need of specifications in the Rust code \cite{deshmukh_prusti_2022}. Furthermore, it can also prove the absence of integer overflows \cite{deshmukh_prusti_2022}. If desired, a developer then can add specifications to the code to get more guarantees \cite{deshmukh_prusti_2022}.

Prusti uses Rust syntax for designing specifications. However, the syntax has been extended by some keywords to help designing those specifications. For example, for creating postconditions, it uses the keyword \textit{result} to refer to the result of a given function \cite{deshmukh_prusti_2022}. Another keyword is \textit{old}, which refers to the state of a given variable, before the function has been executed \cite{deshmukh_prusti_2022}. As can be seen in \autoref{lst:prusti_pre_postcondition}, the Prusti specifications are Rust specific attributes \cite{rust_foundation_attributes_nodate}, distinguished with a Prusti keyword. It is also possible to use pure Rust function calls in the specifications \cite{deshmukh_prusti_2022}. 

\begin{lstlisting}[language=Rust, caption=Example Rust code of a function with a pre- and postcondition for Prusti, label={lst:prusti_pre_postcondition}]
#[requires(*val <= u32::MAX - 1)]
#[ensures(*val == *old(val) + 1)]
fn incr(val: &mut u32) {
    *val += 1;
}
\end{lstlisting}

\subsubsection{Technology}
The backbone of the verification process in Prusti is the Viper Infrastructure, which also has been developed at the ETH Zurich \cite{jobstmann_viper_2016}. The core of this infrastructure is the Viper intermediate language, which can be compared to other intermediate languages like Boogie or Why3 \cite{jobstmann_viper_2016}. The key advantage of the Viper intermediate language is, compared to the other intermediate languages, the support for permission-based reasoning, in particular \textit{implicit dynamic frames} \cite{deshmukh_prusti_2022, astrauskas_leveraging_2019, jobstmann_viper_2016}.

Viper utilizes verification condition generation to Boogie and symbolic execution via Z3 for its verifying backend \cite{jobstmann_viper_2016}. Viper also propagates as one of its three main purposes the interoperability of its verification backend \cite{jobstmann_viper_2016}. For example, Lasse F. Wolff Anthony has shown \cite{anthony_supporting_2022} that it is possible to swap the default Viper SMT solver Z3 with CVC5.

Prusti generates fully automatically so called \textit{core proofs} in Viper, which represents the guarantees the standard rustc compiler already provides, e.g. for borrowing \cite{astrauskas_leveraging_2019}. Those core proofs are necessary, because they provide the foundation for further verification in Viper \cite{astrauskas_leveraging_2019}.

Prusti acts as a frontend for Viper. There are several other Viper frontends for languages like Java, Scala and OpenCL \cite{jobstmann_viper_2016}.

\subsubsection{Feedback}
Due to the tight integration of Prusti in the rustc compiler, verification error messages are displayed like regular Rust compiler error messages. Prusti highlights the relevant specification where the error originates from. However, the user has to figure out by himself what part of the specification exactly is causing the error. To illustrate this, see \autoref{lst:prusti_wrong_postcondition}: the postcondition requires that the function increments the value of a given variable by two, though the function body clearly shows that the function only increments by one. Prusti catches this error and reports that the postcondition might not hold with the given function implementation. The \textit{might} is important here as Prusti is a conservative verifier, which means that it reports always an error, if it is unable to prove that a certain condition holds \cite{prusti_user_guide_absence_2023}.

\begin{lstlisting}[language=bash, caption=Example output of a verification error in Prusti, label={lst:prusti_wrong_postcondition}]
error: [Prusti: verification error] postcondition might not hold.
 --> src/main.rs:4:11
  |
4 | #[ensures(*val == *old(val) + 2)]
  |           ^^^^^^^^^^^^^^^^^^^^^
  |
note: the error originates here
 --> src/main.rs:5:1
  |
5 | / fn incr(val: &mut u32) {
6 | |     *val += 1;
7 | | }
  | |_^

Verification failed
\end{lstlisting}

The Prusti Assistant for Visual Studio Code provides even more convenience, as it highlights the error causing specification directly in the editor.

\subsubsection{Language Support and Verification Conditions}
\label{cap:prusti_language_verification}

Prusti can only prove partial correctness \cite{prusti_user_guide_verification_2023}. The user has to ensure that the Rust code terminates, otherwise the guarantees made by Prusti may not hold up \cite{prusti_user_guide_verification_2023}.

Prusti currently only supports the safe Rust dialect. Support for the unsafe part of Rust is planned \cite{astrauskas_leveraging_2019}, though it is possible today to encapsulate unsafe Rust code and exclude it from the verifying process with the \lstinline{#[trusted]} annotation \cite{prusti_user_guide_trusted_2023}. This is also helpful if a function uses techniques, which aren't supported by Prusti (yet) \cite{prusti_user_guide_trusted_2023}.

Prusti supports designing pre- and postconditions for functions, in a \textit{contract} like manner \cite{deshmukh_prusti_2022, astrauskas_leveraging_2019}. It is possible to call \textit{pure} Rust functions in those conditions, as well as calling \textit{predicates} \cite{prusti_user_guide_predicates_2023}, which have been designed with the \lstinline{predicate!()} macro \cite{deshmukh_prusti_2022}. Those predicates can be equipped with quantifiers \cite{deshmukh_prusti_2022}. 

A key Prusti feature is the support of so called \textit{pledges} \cite{astrauskas_leveraging_2019}. Pledges are used to design specifications for values, which get returned as a reborrow from a function call. With the keywords \textit{before\_expiry} and \textit{after\_expiry}, it is possible to define properties before or after the borrow ends, respectively \cite{prusti_user_guide_pledges_2023}.

Prusti has support for external specifications. This is necessary for function calls to libraries such as the Rust Standard Library, because they don't include any specifications for Prusti. With external specifications, it is possible to extend those function calls with pre- and postconditions and enable Prusti to reason about those calls \cite{deshmukh_prusti_2022,prusti_user_guide_external_2023}. 

Writing specifications for closures is not supported yet by Prusti \cite{prusti_user_guide_closures_2023}. However, Wolff et al. have shown that it is feasible to integrate support for closures into Prusti \cite{wolff_modular_2021}.

Additionally, Prusti has support for
\begin{itemize}
    \item Recursive function calls \cite{deshmukh_prusti_2022}
    \item Inline assertions, refutations and assumptions \cite{prusti_user_guide_assertions_2023}
    \item Type and loop body invariants \cite{deshmukh_prusti_2022}
    \item Type-conditional spec refinements \cite{prusti_user_guide_type-conditional_2023}
    \item Implications \cite{prusti_user_guide_specification_2023}
    \item Snapshot equality \cite{prusti_user_guide_specification_2023}
\end{itemize}

\subsubsection{Development}
Prusti has been developed at the ETH Zurich as a proof-of-concept alongside the publication \textit{Leveraging Rust Types for Modular Specification and Verification} \cite{astrauskas_leveraging_2019}. The Swiss National Science Foundation, Facebook and Amazon Web Services are funding the project \cite{eth_zurich_prusti_2023}. The current projects members from the ETH Zurich are Vytautas Astrauskas, Aurel Bílý, Jonas Fiala, Peter Müller and Federico Poli \cite{eth_zurich_prusti_2023}. Further members are Christoph Matheja from the DTU Copenhagen and Alex Summers from the UBC in Vancouver, Canada \cite{eth_zurich_prusti_2023}.

\subsubsection{Activity}
Prusti is getting actively maintained and improved with over 7,000 commits since 2017 \cite{eth_zurich_viperprojectprusti-dev_2023}. Its latest stable release was in January of 2023, the latest nightly pre-release in May of 2023 \cite{eth_zurich_releases_2023}. 

Prusti is not feature-complete yet; there are about 200 issues in GitHub currently open \cite{eth_zurich_issues_2023}. Many of them are proposals for new features. About 200 issues have been closed already \cite{eth_zurich_issues_2023}.

The Prusti development team has evaluated the tool in several ways: first, they successfully tested the automated core proof generation on 11,791 functions of the top 500 Rust crates of the year 2018 \cite{astrauskas_leveraging_2019}. Second, they identified 519 possible functions in those crates where theoretically an overflow or division by zero could happen, showing that 467 functions have some sort of problem, with the vast majority having an overflow or division by zero problem \cite{astrauskas_leveraging_2019}. Third, they took 16 examples for verifying panic and overflow freedom, while designing and applying specifications on 7 of them \cite{astrauskas_leveraging_2019}. Those 16 examples had a total amount of 471 lines of code, with a total investment of 389 lines of code for Prusti annotations.

Gijsberts has demonstrated that it is feasible to verify a key-value store with Prusti \cite{stef_gijsberts_prusti_2023}.

\subsubsection{Scientific activity}
Prusti is quite popular in the Rust verification community. For example, the tool has been presented at the first occurrence of the Rust Verification Workshop in 2021 \cite{peter_muller_eth_zurich_rust_2021}. This presentation is also featured on the official Rust YouTube channel \cite{alex_summers_prusti_2021}.

The first paper \cite{astrauskas_leveraging_2019} which introduced Prusti has been cited around 150 times according to Google Scholar. Its successor \cite{deshmukh_prusti_2022} has been cited 17 times at the time of writing.

\subsubsection{Documentation}
Prusti has a dedicated user guide to demonstrate its capabilities and provides a guided tour, though incomplete, to teach new users the usage of Prusti \cite{prusti_user_guide_introduction_2023}. Prusti also provides code examples in its two related papers \cite{deshmukh_prusti_2022, astrauskas_leveraging_2019} and its GitHub page \cite{eth_zurich_viperprojectprusti-dev_2023}.

Prusti ships with over 300 test cases \cite{astrauskas_leveraging_2019}. They include correct and incorrect examples \cite{eth_zurich_prusti-devprusti-tests_2023} of Rust code and annotations.